\documentclass[11pt,a4paper]{article}

\usepackage[margin=2.6cm]{geometry}
\usepackage{amsmath,amssymb,amsthm,mathtools}
\usepackage{bm}
\usepackage{enumitem}
\usepackage{xcolor}
\usepackage[colorlinks=true,linkcolor=blue!50!black,citecolor=blue!50!black]{hyperref}

\theoremstyle{plain}
\newtheorem{theorem}{Theorem}[section]
\newtheorem{lemma}[theorem]{Lemma}
\newtheorem{proposition}[theorem]{Proposition}
\newtheorem{corollary}[theorem]{Corollary}
\theoremstyle{definition}
\newtheorem{definition}[theorem]{Definition}
\newtheorem{assumption}{Assumption}
\renewcommand{\theassumption}{\Alph{assumption}}
\theoremstyle{remark}
\newtheorem{remark}[theorem]{Remark}
\newtheorem{example}[theorem]{Example}

\DeclareMathOperator*{\argmax}{arg\,max}
\DeclareMathOperator*{\argmin}{arg\,min}
\DeclareMathOperator{\dom}{dom}
\newcommand{\R}{\mathbb{R}}
\newcommand{\E}{\mathbb{E}}
\newcommand{\KL}{\mathrm{KL}}
\newcommand{\Expl}{\mathrm{Expl}}
\newcommand{\QRE}{\mathrm{QRE}}
\newcommand{\simplex}{\Delta}
\newcommand{\eps}{\varepsilon}
\newcommand{\half}{\tfrac{1}{2}}
\newcommand{\ip}[2]{\langle #1, #2\rangle}
\newcommand{\norm}[1]{\lVert #1 \rVert}
\newcommand{\abs}[1]{\lvert #1 \rvert}

\title{\textbf{A Simple Convergence Theory for MMD on Fixed-Size Mixtures}\\[4pt]
\large Tabular, exact-gradient, purely about where the gradient pushes}
\author{theory3}
\date{}

\begin{document}
\maketitle

\begin{abstract}
This note gives the simplest convergence argument I can find for the algorithm in
this codebase (magnetic mirror descent / MMD on the categorical head of a
$K$-component mixture policy, plus entropy-regularised ascent on the component
locations), in the tabular, exact-gradient regime.

Everything rests on one exact identity: the exploitability of a pair of mixture
policies splits, with no slack, into the exploitability of a \emph{finite matrix
game} played on the current component locations plus two \emph{support gaps}.
The two halves are then killed by two different mechanisms: the matrix-game half
by MMD on the simplex (which is a contraction because of the magnet), the support
gaps by the component locations climbing the payoff landscape. The two mechanisms
interact, and a two-line small-gain calculation says exactly when the interaction
is benign.

The Gaussian head never appears in the analysis except as a positive scalar
multiplying a gradient (which cannot move a fixed point) and as the smoothing that
defines \emph{which} landscape the gradient sees. Approximation error, sampling
noise, and critic error are assumed away.
\end{abstract}

\tableofcontents

\newpage
\section{What this document is, and how it differs from \texttt{theory2}}

\texttt{theory2} proves a lot, and the proofs are involved: they carry
Wasserstein geometry on measures, heat-kernel annealing of the Gaussian width,
lifted monotonicity, and a five-part incremental construction. This note throws
almost all of that away and asks the narrow question:

\begin{quote}
\emph{In the tabular, noise-free, exact-gradient regime, with a fixed number $K$ of
components and $K$ larger than what is needed to represent a Nash equilibrium ---
what does the gradient field do, and when does it converge?}
\end{quote}

The four simplifications that make it easy:

\begin{enumerate}[label=(S\arabic*),leftmargin=*]
\item \textbf{Exact gradients.} No Monte-Carlo, no critic, no PPO surrogate. The
update is the deterministic vector field of \texttt{idealized\_mmd.py}.
\item \textbf{Atoms, not measures.} A $K$-component mixture is a point in
$\simplex_K \times A^K$: $K$ weights and $K$ locations. We never work in the space
of measures, never need Wasserstein gradient flows, and never need a
representation-error term --- because of (S4).
\item \textbf{The Gaussian head is not the subject.} A Gaussian component
$\mathcal N(x_k,\sigma_k^2)$ is how the algorithm \emph{obtains} a gradient
(it samples in a neighbourhood, so it feels the smoothed landscape). What we
analyse is the vector field that comes out. The width $\sigma$ enters this
document in exactly two harmless places: it multiplies the gradient by a positive
scalar (Lemma~\ref{lem:precond}: cannot move a fixed point), and it fixes which
smoothed landscape $u_\sigma$ the atoms are climbing (\S\ref{sec:setting},
Remark~\ref{rem:sigma}). It is never a state variable of the proof.
\item \textbf{Capacity is not the bottleneck.} We assume $K$ is at least the
number of ``cells'' of the landscape (Assumption~\ref{as:cells}), which is at
least the size of the support of a Nash. So there is no representation bias to
bound and no $\Theta(1/t)$ bias-annealing story.
\end{enumerate}

What survives from \texttt{theory2}, in a much shorter form, is the shape of the
answer: convergence is governed by \emph{coupling strength versus landscape
curvature}, and the magnet temperature is the knob that buys stability at the
price of bias. Theorem~\ref{thm:main} states that as a single scalar inequality,
$\kappa\lesssim m_\star\tau$, and it is proved the way the fixed-matrix MMD result is
proved: one map, one fixed point, one contraction.

\medskip
\noindent\textbf{Notation summary.}
\begin{center}
\begin{tabular}{ll}
$A,B\subset\R^d$ & compact action sets of the maximiser / minimiser\\
$u(a,b)$ & payoff to the maximiser; $\sigma$-smoothed version $u_\sigma$\\
$x_1,\dots,x_K\in A$, $w\in\simplex_K$ & maximiser's atoms and weights, $\mu=\sum_k w_k\delta_{x_k}$\\
$y_1,\dots,y_K\in B$, $v\in\simplex_K$ & minimiser's atoms and weights, $\nu=\sum_l v_l\delta_{y_l}$\\
$M(X,Y)_{kl}=u(x_k,y_l)$ & the \emph{restricted game}: a $K\times K$ matrix\\
$\tau$ & magnet / entropy temperature\\
$\eta_x,\eta_w$ & step sizes for locations and for weights\\
$m$ & strong concavity of the landscape inside a cell\\
$\kappa$ & coupling strength, $\kappa=\sup\norm{\nabla_a u(a,b)-\nabla_a u(a,b')}$\\
\end{tabular}
\end{center}

\newpage
\part*{Part I --- The setting and the identity}
\addcontentsline{toc}{part}{Part I --- The setting and the identity}

\section{Setting}\label{sec:setting}

\paragraph{The game.} A two-player zero-sum game with continuous actions:
compact convex sets $A,B\subset\R^d$ and a payoff $u:A\times B\to\R$ that the
first player maximises and the second minimises. Mixed strategies are probability
measures; the payoff extends bilinearly,
$U(\mu,\nu)=\int\!\!\int u\,d\mu\,d\nu$. Write $u(a,\nu)=\int u(a,b)\,d\nu(b)$
and $u(\mu,b)=\int u(a,b)\,d\mu(a)$.

\paragraph{Exploitability.} The quantity we drive to zero:
\begin{equation}\label{eq:expl}
\Expl(\mu,\nu) \;=\; \max_{a\in A} u(a,\nu)\;-\;\min_{b\in B} u(\mu,b)\;\ge\;0,
\end{equation}
which is $0$ exactly at a Nash equilibrium. (This is \texttt{exploitability()} in
\texttt{idealized\_mmd.py}, computed there on a fine grid.)

\paragraph{The policy class.} Both players use a mixture of $K$ atoms:
\[
\mu=\sum_{k=1}^K w_k\,\delta_{x_k},\qquad
\nu=\sum_{l=1}^K v_l\,\delta_{y_l},\qquad w,v\in\simplex_K,\; x_k\in A,\; y_l\in B .
\]
The state of the algorithm is $z=(X,w,Y,v)\in A^K\times\simplex_K\times B^K\times\simplex_K$.

\paragraph{Standing regularity.} $u$ is $C^2$ with
\[
\abs{u}\le U_{\max},\qquad
\norm{\nabla_a u}\le G,\qquad
\norm{\nabla^2_{aa} u}\le L,\qquad
\kappa:=\sup_{a,b,b'}\norm{\nabla_a u(a,b)-\nabla_a u(a,b')}\le 2G ,
\]
and symmetrically in $b$. The constant $\kappa$ is the one that matters: it is
\emph{how much the maximiser's landscape moves when the minimiser moves}. In a game
$u(a,b)=D(a)-D(b)+\lambda\,f(a)^\top g(b)$ --- the family in \texttt{games/examples.py}
--- we have $\kappa\le 2\lambda\,\norm{f'}_\infty\norm{g}_\infty$, i.e. $\kappa$ is
literally the \texttt{coupling} coefficient up to a constant. If $\kappa=0$ the two
players play separate single-agent problems and everything below is trivial.

\begin{remark}[Where the Gaussian width lives]\label{rem:sigma}
The algorithm's components are $\mathcal N(x_k,\sigma_k^2 I)$, not atoms, so the
gradient it computes is a gradient of the \emph{smoothed} payoff
$u_\sigma(a,b) = \E_{\xi\sim\mathcal N(0,\sigma^2 I)}[u(a+\xi,b)]$
(this is exactly what \texttt{well\_expectation} computes in closed form).
Fix $\sigma$ in a compact range $[\sigma_{\min},\sigma_{\max}]$ --- the code clips
$\log\sigma$ to $[\log 10^{-3},0]$ --- and read \emph{every} $u$ in this document as
$u_\sigma$. All constants $G,L,\kappa,m$ are then constants of $u_\sigma$. Nothing else
changes; in particular Assumption~\ref{as:cells} is an assumption about $u_\sigma$,
not about $u$. This is the only role $\sigma$ plays, and \S\ref{sec:anneal-sigma}
says what it costs to let $\sigma$ move.
\end{remark}

\section{The algorithm, tabularised}\label{sec:alg}

The implementation is soft actor-critic with a magnetic-mirror-descent actor. In
the tabular exact-gradient regime the critic is exact, so the whole thing collapses
to two updates. Write $q_k = u(x_k,\nu_t)$ for the maximiser's per-component value
(this is \texttt{component\_q}); note $\ip{w}{q} = U(\mu,\nu)$.

\paragraph{(W) The weights: closed-form MMD on the simplex.} With magnet $\rho$,
magnet temperature $\tau$, entropy temperature $\tau_{\mathrm{ent}}$ and step $\eta_w$,
\begin{equation}\label{eq:mmd}
w_{t+1} \;=\; \argmax_{w\in\simplex_K}\;
\Big\{ \ip{w}{q} - \tau\KL(w\Vert\rho) - \tau_{\mathrm{ent}}\KL(w\Vert \mathrm{unif})
- \tfrac{1}{\eta_w}\KL(w\Vert w_t)\Big\},
\end{equation}
whose solution is the log-linear update implemented in
\texttt{categorical\_mirror\_update}:
$\log w_{t+1} \propto (\eta_w q + \eta_w\tau\log\rho + \log w_t)/(1+\eta_w\tau+\eta_w\tau_{\mathrm{ent}})$.
The magnet $\rho$ is refreshed to the current $w$ every $T_{\mathrm{mag}}$ steps.
Since $\KL(\cdot\Vert\mathrm{unif})$ is just $\KL$ against a particular magnet, we
absorb $\tau_{\mathrm{ent}}$ into $\tau$ from now on and write a single temperature
$\tau>0$. \emph{SAC's entropy bonus and MMD's magnet are the same object here.}

\paragraph{(X) The locations: preconditioned ascent.} The mixture's payoff is
$U(\mu,\nu)=\sum_k w_k\,u(x_k,\nu)$, so
$\nabla_{x_k} U = w_k\,\nabla_a u(x_k,\nu)$; the natural gradient in the Gaussian's
Fisher metric multiplies this by $\sigma_k^2$. Hence the implemented update
(\texttt{gaussian\_natural\_step}) is, exactly,
\begin{equation}\label{eq:xstep}
x_k \;\leftarrow\; \Pi_A\big(x_k + \eta_x\,p_k\,\nabla_a u(x_k,\nu_t)\big),
\qquad p_k \;=\; \sigma_k^2\,w_k \;>\;0 .
\end{equation}
The minimiser's updates are \eqref{eq:mmd}--\eqref{eq:xstep} with $u$ replaced by
$-u$.

\begin{remark}[What \eqref{eq:xstep} leaves out]\label{rem:xmagnet}
\texttt{gaussian\_natural\_step} actually ascends
$U(\mu,\nu)+\tau_{\mathrm{ent}}\sum_k\log\sigma_k
-\tau\,\KL\big(\mathcal N(x_k,\sigma_k^2)\Vert\mathcal N(\rho_{X,k},\sigma_{\rho,k}^2)\big)$,
so the implemented mean update carries, on top of \eqref{eq:xstep}, a proximal pull
$-\eta_x\sigma_k^2\,\tau\,(x_k-\rho_{X,k})/\sigma_{\rho,k}^2$ towards the magnet's
means. Everything below holds verbatim with that term included: it only \emph{adds}
$\tau\sigma_k^2/\sigma_{\rho,k}^2$ to the strong-concavity constant $m$ of the
location objective, and it moves the cell maximiser to that of the proximal
objective. It vanishes at every magnet refresh $\rho_X\leftarrow X$, so it does not
change the outer fixed point. We drop it because it can only help.
\end{remark}

\paragraph{(\texorpdfstring{$\sigma$}{sigma}) The widths.} Not analysed; see
Remark~\ref{rem:sigma}. We only use $p_k\in[p_{\min},p_{\max}]$ with
$p_{\min}=\sigma_{\min}^2 w_{\min}>0$; \S\ref{sec:starve} is about keeping
$w_{\min}$ away from $0$.

\begin{lemma}[Preconditioner neutrality]\label{lem:precond}
Let $A=\prod_{i=1}^d[\ell_i,r_i]$ be a box and $P=\mathrm{diag}(p_1,\dots,p_d)$ with
every $p_i>0$. Then the fixed points of $x\mapsto \Pi_A(x+\eta P\nabla\phi(x))$ are
exactly the stationary points of $\phi$ on $A$, independently of $P$.
\end{lemma}
\begin{proof}
$\Pi_A$ is the coordinatewise clip, so with $g=\nabla\phi(x)$ we have
$\Pi_A(x+\eta Pg)=x$ iff, for every coordinate $i$, either $p_ig_i=0$, or
$x_i=r_i$ and $p_ig_i>0$, or $x_i=\ell_i$ and $p_ig_i<0$. Each clause is unchanged
when $p_ig_i$ is replaced by $g_i$, because $p_i>0$; and the resulting condition is
exactly $g\in N_A(x)$, i.e.\ stationarity of $\phi$ on $A$.
\end{proof}

\begin{remark}
$P\succ0$ alone does \emph{not} suffice on a general convex $A$: being fixed asks
$Pg\in N_A(x)$, and a non-diagonal $P$ can rotate $g$ into or out of the normal cone
at a curved boundary. A box with diagonal $P$ --- the case in this codebase, where
$A$ is an interval and $P$ is a per-component scalar --- is exactly the situation in
which $N_A(x)$ is $P$-invariant.
\end{remark}

\noindent This one line is the formal content of ``the analysis should not be
about the Gaussian head''. The head chooses $P$; $P$ changes speed, never
destination. It \emph{can} change stability in a genuinely coupled two-player
system. It is harmless here for a narrower reason: the location block is projected
gradient ascent--descent on a single \emph{strongly monotone} saddle problem
(\S\ref{sec:locations}), whose contraction survives any constant positive diagonal
preconditioner once the metric is chosen to match it (Remark~\ref{rem:Pnorm}). That is
why the final condition (Theorem~\ref{thm:main}) depends on $p_{\max}/p_{\min}$ only as
a constant.

\section{The master identity}\label{sec:identity}

Everything in this note is a consequence of the following exact decomposition.
Define, for the current state $z=(X,w,Y,v)$:
\begin{align}
\Delta_x(z) &:= \max_{a\in A} u(a,\nu) \;-\; \max_{k\le K} u(x_k,\nu)
&&\text{(maximiser's \emph{support gap})}\\
\Delta_y(z) &:= \min_{l\le K} u(\mu,y_l) \;-\; \min_{b\in B} u(\mu,b)
&&\text{(minimiser's support gap)}\\
\Expl_M(w,v) &:= \max_{k} (Mv)_k \;-\; \min_{l} (w^\top M)_l,
\quad M=M(X,Y) &&\text{(\emph{matrix-game} exploitability).}
\end{align}
All three are $\ge0$: the first two because a max over a subset is at most a max
over the whole set, the third because $\max_k(Mv)_k\ge w^\top Mv\ge\min_l(w^\top M)_l$.

\begin{lemma}[Master identity]\label{lem:identity}
For every state $z$,
\[
\boxed{\;\Expl(\mu,\nu)\;=\;\Expl_{M(X,Y)}(w,v)\;+\;\Delta_x(z)\;+\;\Delta_y(z).\;}
\]
\end{lemma}
\begin{proof}
Insert a telescoping pair into \eqref{eq:expl}:
\[
\max_a u(a,\nu)-\min_b u(\mu,b)
=\underbrace{\Big[\max_a u(a,\nu)-\max_k u(x_k,\nu)\Big]}_{\Delta_x}
+\Big[\max_k u(x_k,\nu)-\min_l u(\mu,y_l)\Big]
+\underbrace{\Big[\min_l u(\mu,y_l)-\min_b u(\mu,b)\Big]}_{\Delta_y}.
\]
Finally $u(x_k,\nu)=\sum_l v_l u(x_k,y_l)=(Mv)_k$ and
$u(\mu,y_l)=\sum_k w_k u(x_k,y_l)=(w^\top M)_l$, so the middle bracket is
$\Expl_M(w,v)$.
\end{proof}

\begin{corollary}[Support sufficiency]\label{cor:suff}
If (i) some atom of each player is a global best response, i.e.
$\Delta_x=\Delta_y=0$, and (ii) $(w,v)$ is a Nash equilibrium of the $K\times K$
matrix game $M(X,Y)$, then $(\mu,\nu)$ is a Nash equilibrium of the continuous
game.
\end{corollary}

\medskip
Lemma~\ref{lem:identity} is the whole architecture of the theory, and it is worth
saying plainly what it buys:

\begin{itemize}[leftmargin=*]
\item \textbf{The hard part of game theory --- mixing --- is finite.} Once the
locations are fixed, finding the right \emph{weights} is a $K\times K$ matrix game
problem, and MMD solves matrix games. Nothing about continuous action spaces,
non-concavity, or infinite-dimensional geometry enters that half.
\item \textbf{The hard part of continuous actions --- \emph{where} to put mass ---
carries no equilibrium structure of its own.} $\Delta_x$ is a functional of $(X,\nu)$
alone: no $\max\min$, no equilibrium concept appears in it, in contrast with $\Expl_M$,
which is a genuine saddle quantity. Correspondingly, the update that closes it
consults no equilibrium either --- it is ascent on a function.

\emph{This must not be over-read.} It does \emph{not} say that closing $\Delta_x$ is a
single-agent problem. At $K=1$ the matrix $M$ is $1\times1$, so $\Expl_M\equiv0$ and the
identity degenerates to $\Expl=\Delta_x+\Delta_y$ --- the whole of the game now sits in
the ``non-strategic'' half, and the dynamics are simultaneous gradient
ascent--descent on $u(x,y)$, which on $u=xy$ cycles forever. The support gaps are closed
against an opponent who is himself moving, and moving in response to us; that is a game,
with feedback, and \S\ref{sec:main} is the price. What makes the location half tractable
is not the identity but Assumption~\ref{as:cells}(C1): strong concavity per cell makes
the joint location dynamics a \emph{strongly monotone} game, and strongly monotone games
behave like optimisation problems. The honest reading of the split is therefore
\emph{strategic error versus support error}, not \emph{game versus optimisation}.
\item \textbf{They are added, not multiplied.} There is no cross term. So the two
error sources can be attacked by two independent mechanisms, and the only thing
left to control is how each mechanism disturbs the other. That is \S\ref{sec:main}.
\end{itemize}

\newpage
\part*{Part II --- The two mechanisms}
\addcontentsline{toc}{part}{Part II --- The two mechanisms}

\section{The one assumption about the game: cells}\label{sec:cells}

$\Delta_x$ is closed by gradient ascent, and gradient ascent is local. So we need
an assumption saying that local ascent from a spread-out initialisation finds the
global maximum. The weakest useful version:

\begin{assumption}[Cell decomposition]\label{as:cells}
There is a finite family of compact \emph{convex} sets $A_1,\dots,A_{m_A}\subseteq A$
and a constant $m>0$ such that for \emph{every} $\nu\in\simplex(B)$:
\begin{enumerate}[label=(C\arabic*),leftmargin=*]
\item \textbf{Curvature.} $a\mapsto u(a,\nu)$ is $m$-strongly concave (and $L$-smooth)
on each $A_i$.
\item \textbf{Invariance.} Each $A_i$ is invariant under the update map
$a\mapsto\Pi_A\big(a+\eta_x P\nabla_a u(a,\nu)\big)$, for every admissible diagonal
$P$ and every $\eta_x$ obeying the step-size condition of Lemma~\ref{lem:track}.
\item \textbf{Reachability of the best response.} Some cell contains a global
maximiser of $u(\cdot,\nu)$ over $A$.
\end{enumerate}
Symmetrically for $B$ with cells $B_1,\dots,B_{m_B}$ and the same $m$.
\end{assumption}

By (C1) each cell has a unique maximiser
$x_i^\star(\nu) := \argmax_{a\in A_i} u(a,\nu)$, and (C3) says exactly
\begin{equation}\label{eq:cellmax}
\max_{a\in A} u(a,\nu) = \max_{i\le m_A} u\big(x_i^\star(\nu),\nu\big).
\end{equation}
So the global maximum \emph{is} the best of the per-cell maxima, and each per-cell
maximum is reachable by local ascent. That is the entire content of the assumption.

\begin{remark}[Why the cells must \emph{not} be asked to cover $A$]\label{rem:nocover}
The tempting form of (C3) is $\bigcup_iA_i=A$, and it is inconsistent with (C1).
If $u(\cdot,\nu)$ has an interior local \emph{minimum} $a_0$ --- which every
multi-well landscape has, and \texttt{MultiPointGame} certainly has --- then $a_0$
would lie in some cell on which $u(\cdot,\nu)$ is $m$-strongly concave, contradicting
$\nabla^2_{aa}u(a_0,\nu)\succeq0$. (C3) is all that \eqref{eq:cellmax} ever needs, and
it is consistent: cells are neighbourhoods of the peaks, the valleys and plateaus
between them belong to no cell, and \S\ref{sec:anneal-sigma} is precisely about making
the peaks wide enough that the initialisation lands inside one.
\end{remark}

\begin{remark}[Invariance is stated for the \emph{discrete} step]
(C2) is about the map the algorithm runs, not about the flow $\dot a=\nabla_au(a,\nu)$,
and the algorithm projects onto $A$, not onto $A_i$; a continuous-time invariance
statement would not by itself keep a discrete iterate inside its cell. A sufficient
condition is that $A_i$ contain the $\norm{\cdot}_{P^{-1}}$-ball of radius
$\sup_\nu\mathrm{diam}(A_i)$ around $x_i^\star(\nu)$ for all $\nu$, since by
Lemma~\ref{lem:track} the step is a contraction towards $x_i^\star(\nu)$ in that norm.
\end{remark}

\begin{assumption}[Capacity and covering initialisation]\label{as:cover}
$K\ge m_A$ and $K\ge m_B$, and the initial atoms $x_1(0),\dots,x_K(0)$ meet every
cell $A_i$ (symmetrically for $Y$). Write $c(k)$ for a cell containing $x_k(0)$; by
(C2), $x_k(t)\in A_{c(k)}$ for every $t$.
\end{assumption}

This is the formal reading of ``we always have more points than we need to
reconstruct a Nash''. It is genuinely stronger than that phrase: the support of a
Nash is contained in the set of best responses, which by \eqref{eq:cellmax} contains
at most one point per cell, so $\lvert\mathrm{supp}\,\mu^\star\rvert \le m_A \le K$
automatically. Assumption~\ref{as:cover} additionally asks that the atoms be
\emph{spread across} the cells at initialisation, because ascent will never move an
atom between cells. In the code this is
\texttt{\_spread\_bias\_init} / the spread initialisations in
\texttt{run\_idealized.py}, and \S\ref{sec:anneal-sigma} explains the other way to
get it (start with a large $\sigma$, which merges cells).

\begin{lemma}[How far a mixture can move a gradient]\label{lem:grad}
For mixtures $\nu=\sum_lv_l\delta_{y_l}$ and $\nu'=\sum_lv'_l\delta_{y'_l}$ and any $a$,
\[
\norm{\nabla_a u(a,\nu)-\nabla_a u(a,\nu')}\;\le\;\tfrac{\kappa}{2}\norm{v-v'}_1
\;+\;L_{ab}\max_l\norm{y_l-y'_l},\qquad L_{ab}:=\norm{\nabla^2_{ab}u}_\infty .
\]
\end{lemma}
\begin{proof}
Split
$\nabla_au(a,\nu)-\nabla_au(a,\nu')=\sum_l(v_l-v'_l)\nabla_au(a,y_l)
+\sum_lv'_l\big(\nabla_au(a,y_l)-\nabla_au(a,y'_l)\big)$.
The second sum is at most $L_{ab}\max_l\norm{y_l-y'_l}$. The first has coefficients
summing to $0$, so subtracting any fixed vector from every $\nabla_au(a,y_l)$ leaves it
unchanged; subtracting the midpoint of their range and using the definition of $\kappa$
bounds it by $\tfrac{\kappa}{2}\norm{v-v'}_1$.
\end{proof}

\noindent The ratio $\kappa/m_\star$ --- \textbf{coupling over curvature} --- is the
single number that decides whether this algorithm is stable. It appears again in
Theorem~\ref{thm:main} and nowhere else. Note which constant it is \emph{not}: the
cross-curvature $L_{ab}=\norm{\nabla^2_{ab}u}_\infty$, the second term above, never
reaches the threshold --- see Remark~\ref{rem:noLab}.

\section{The weight loop is a contraction (because of the magnet)}\label{sec:weights}

Fix the locations, so the players face the fixed matrix game $M$. With temperature
$\tau>0$ and magnet $\rho$, the regularised saddle problem
\begin{equation}\label{eq:qre}
\max_{w\in\simplex_K}\min_{v\in\simplex_K}\;
\Phi_M(w,v),\qquad
\Phi_M(w,v):=w^\top M v - \tau\KL(w\Vert\rho_x) + \tau\KL(v\Vert\rho_y)
\end{equation}
is $\tau$-strongly-convex-concave (KL is $1$-strongly convex in $\norm\cdot_1$), so it
has a unique solution $\hat z(M)=(\hat w,\hat v)$, the \emph{quantal response
equilibrium} $\QRE_\tau(M)$. Two standard facts, both of which we import:

\begin{proposition}[MMD contracts, Sokota et al.\ 2023]\label{prop:mmd}
Write $z_W=(w,v)$ and let
\[
D(z'_W\Vert z_W)\;:=\;\KL(w'\Vert w)+\KL(v'\Vert v)
\]
be the Bregman divergence of negative entropy on $\simplex_K\times\simplex_K$. For a
fixed $M$ with $\norm{M}_{\max}\le U_{\max}$ and $\eta_w\le\eta_0(\tau,U_{\max})$, one
simultaneous MMD step \eqref{eq:mmd} for both players, $z_W\mapsto F_M(z_W)$, satisfies
\[
D\big(\hat z(M)\,\big\Vert\,F_M(z_W)\big)\;\le\;(1-c\,\eta_w\tau)\,
D\big(\hat z(M)\,\big\Vert\,z_W\big)
\]
for an absolute constant $c>0$.
\end{proposition}

\begin{remark}[The divergence is not decoration]
It is tempting to state Proposition~\ref{prop:mmd} as an $\ell_1$ contraction
$\norm{F_M(z)-\hat z}_1\le(1-c\eta_w\tau)\norm{z-\hat z}_1$, and tempting further to
state it as a \emph{two-point} contraction $\norm{F_M(z)-F_M(z')}\le(1-c\eta_w\tau)\norm{z-z'}$,
which would let one apply Banach's theorem directly. Both are false in general. The
argument behind Proposition~\ref{prop:mmd} is a Lyapunov argument: the bilinear part of
the saddle operator cancels in $\ip{F_M(z)-F_M(\hat z)}{z-\hat z}$ by monotonicity, and
only the KL part survives. A two-point bound destroys that cancellation --- in logit
coordinates the simultaneous step expands by $(1+\eta_wU_{\max})/(1+\eta_w\tau)$, which
exceeds $1$ whenever $U_{\max}>\tau$. So the contraction really is towards the fixed
point, in $D$, and the rest of this note is organised around that.
\end{remark}

\begin{lemma}[One MMD step taken against the wrong matrix]\label{lem:mmdpert}
Let $\bar M$ be any payoff matrix, $\bar z_W=\hat z(\bar M)$ its QRE, and let
$z_W^+=F_M(z_W)$ be one MMD step taken with a possibly different matrix $M$. Then
\[
D(\bar z_W\Vert z_W^+)\;\le\;(1-c\,\eta_w\tau)\,D(\bar z_W\Vert z_W)
\;+\;2\,\eta_w\,\norm{M-\bar M}_{\max}\,\norm{z_W-\bar z_W}_1 .
\]
\end{lemma}
\begin{proof}
Let
$\mathcal F_M(z_W)=\big(-Mv+\tau\nabla\KL(w\Vert\rho_x),\;M^\top w+\tau\nabla\KL(v\Vert\rho_y)\big)$
be the saddle operator of \eqref{eq:qre}; the closed form in \eqref{eq:mmd} is the mirror
step $z_W^+=\argmin_{z'}\{\eta_w\ip{\mathcal F_M(z_W)}{z'}+D(z'\Vert z_W)\}$ with the KL
terms taken proximally. Proposition~\ref{prop:mmd} is exactly the conclusion of the
three-point mirror-descent inequality for $M=\bar M$, where $\tau$-strong monotonicity
of $\mathcal F_{\bar M}$ supplies the factor $(1-c\eta_w\tau)$. Running the identical
argument with the operator $\mathcal F_M$ while comparing against $\bar z_W$ changes one
term: it adds $\eta_w\ip{\mathcal F_M(z_W)-\mathcal F_{\bar M}(z_W)}{z_W-\bar z_W}$. Now
$\mathcal F_M(z_W)-\mathcal F_{\bar M}(z_W)=\big(-(M-\bar M)v,\;(M-\bar M)^\top w\big)$,
whose two blocks have $\ell_\infty$-norm at most $\norm{M-\bar M}_{\max}$ each, so
H\"older bounds the added term by $2\eta_w\norm{M-\bar M}_{\max}\norm{z_W-\bar z_W}_1$.
\end{proof}

\noindent Lemma~\ref{lem:mmdpert} is what replaces a sensitivity analysis of the QRE.
We never need to know how far $\hat z(M)$ moves when $M$ moves; we only need that one
step taken with the wrong matrix still contracts towards the \emph{right} fixed point,
up to a disturbance linear in the matrix error. That is the single structural change
that makes the rest of this note shorter --- and, as \S\ref{sec:main} shows, sharper ---
than the version that tracked moving targets.

\begin{lemma}[Bias of the QRE]\label{lem:bias}
$\Expl_M(\hat w,\hat v)\;\le\;2\tau\log(1/\rho_{\min})$, where
$\rho_{\min}=\min_k\min(\rho_{x,k},\rho_{y,k})$. In particular with a uniform
magnet, $\Expl_M(\QRE_\tau)\le 2\tau\log K$.
\end{lemma}
\begin{proof}
Standard: $(\hat w,\hat v)$ is an exact Nash of the perturbed game with payoff
$w^\top Mv + \tau\KL(w\Vert\rho_x)-\tau\KL(v\Vert\rho_y)$, and each perturbation term
ranges over an interval of length at most $\tau\log(1/\rho_{\min})$ on $\simplex_K$.
\end{proof}

\begin{lemma}[No atom starves]\label{lem:full}
\emph{(i)} $\hat w_k \;\ge\; \rho_{x,k}\,e^{-2U_{\max}/\tau}$ for all $k$.
\emph{(ii)} The same holds along the iterates, not merely at the fixed point: with
$\psi_t:=\max_k\log w_{t,k}-\min_k\log w_{t,k}$ and $\psi_\rho$ the same quantity for
$\rho_x$,
\[
\psi_t\;\le\;\max\{\psi_0,\;\psi_\rho+2U_{\max}/\tau\}\quad\text{for all }t,
\qquad\text{hence}\qquad w_{\min}(t)\;\ge\;K^{-1}e^{-\psi_t}\;>\;0 .
\]
\end{lemma}
\begin{proof}
(i) The first-order condition of \eqref{eq:qre} is
$\hat w \propto \rho_x\exp\big((M\hat v)/\tau\big)$, and $\norm{M\hat v}_\infty\le U_{\max}$,
so the exponential factor lies in $[e^{-U_{\max}/\tau},e^{U_{\max}/\tau}]$ and the
normaliser in the reciprocal range.

(ii) Put $\theta=1/(1+\eta_w\tau)$. The closed form in \eqref{eq:mmd} gives
$\log w_{t+1,k}=\theta\log w_{t,k}+(1-\theta)\log\rho_{x,k}+\theta\eta_wq_k-\log Z_t$,
so for any $k,j$ the log-ratio obeys
$\log\frac{w_{t+1,k}}{w_{t+1,j}}=\theta\log\frac{w_{t,k}}{w_{t,j}}
+(1-\theta)\log\frac{\rho_{x,k}}{\rho_{x,j}}+\theta\eta_w(q_k-q_j)$,
whence $\psi_{t+1}\le\theta\psi_t+(1-\theta)\psi_\rho+2\theta\eta_wU_{\max}$. Since
$(1-\theta)=\eta_w\tau\theta$, the fixed point of this scalar recursion is
$\psi_\rho+2U_{\max}/\tau$, and the recursion is a contraction, so $\psi_t$ never
exceeds the maximum of $\psi_0$ and that value. Finally $\max_kw_{t,k}\ge1/K$.
\end{proof}

\noindent Part (ii) matters: the location step \eqref{eq:xstep} is scaled by the
\emph{current} weight $w_k(t)$, not by $\hat w_k$, so a bound that holds only at the
QRE would be circular. Part (ii) is an invariant of the update and needs nothing from
the rest of the theory.

\noindent Lemma~\ref{lem:full} is the quiet hero of the theory. The regularisation
that makes the weight loop a contraction \emph{also} guarantees that every atom keeps
a strictly positive weight, and therefore --- because the location step
\eqref{eq:xstep} is scaled by $w_k$ --- that every atom keeps moving. Unregularised
mixture ascent kills its own exploration: an atom that temporarily looks bad gets
weight $\approx0$, freezes, and can never be revived. See \S\ref{sec:starve}.

\section{The location loop is a contraction (inside a cell)}\label{sec:locations}

\paragraph{The location block is one saddle problem, not two ascent problems.}
It is tempting --- and it is what an earlier draft of this section did --- to analyse
each player's atoms separately, as ascent on $a\mapsto u(a,\nu_t)$ with the opponent's
atoms treated as an exogenous disturbance. That framing is available, but it is lossy:
it forces an extra hypothesis bounding the cross-curvature $\norm{\nabla^2_{ab}u}$
against $m$, which the algorithm does not actually need. The reason is that the two
location blocks are \emph{not} two disturbances to each other; together they are
projected gradient ascent--descent on a single saddle function, and the
$\nabla^2_{ab}u$ terms cancel between the blocks exactly as they do in any saddle
problem. Write, with the weights held at their current values,
\begin{equation}\label{eq:Udef}
U(X,Y)\;=\;\sum_{k,l}w_kv_l\,u(x_k,y_l),\qquad
\nabla_{x_k}U \;=\; w_k\,\nabla_a u(x_k,\nu),\qquad
\nabla_{y_l}U \;=\; v_l\,\nabla_b u(\mu,y_l).
\end{equation}
The implemented step \eqref{eq:xstep}, $x_k\leftarrow\Pi_A(x_k+\eta_x\sigma_k^2w_k\nabla_au(x_k,\nu))$,
is therefore \emph{exactly}
\[
x_k\;\leftarrow\;\Pi_A\big(x_k+\eta_x\,\sigma_k^2\,\nabla_{x_k}U\big),
\]
i.e.\ preconditioned ascent on $U$ with the \textbf{constant} preconditioner
$P_0=\mathrm{diag}(\sigma_k^2)$. The factor $w_k$ that \eqref{eq:xstep} shows in the
preconditioner belongs, more usefully, to the gradient. Two things follow at once: the
metric no longer moves with $t$ (Remark~\ref{rem:Pnorm}), and the joint update is
ascent--descent on one function.

\medskip
\noindent Throughout, $P_0=\mathrm{diag}(p_1,\dots)$ is constant positive diagonal with
$p_{\min}I\preceq P_0\preceq p_{\max}I$ ($p_{\min}=\sigma_{\min}^2$,
$p_{\max}=\sigma_{\max}^2$), and $\norm{\zeta}_{P_0^{-1}}^2:=\zeta^\top P_0^{-1}\zeta$, so
$p_{\min}\norm{\zeta}^2_{P_0^{-1}}\le\norm{\zeta}^2\le p_{\max}\norm{\zeta}^2_{P_0^{-1}}$.
Since $A,B$ are boxes and $P_0$ is diagonal, the projection is a coordinatewise clip and
is nonexpansive in $\norm{\cdot}_{P_0^{-1}}$. Write $\xi=(X,Y)$ for the joint location
vector, $\mathcal C=\prod_kA_{c(k)}\times\prod_lB_{c(l)}$ for its (convex, compact)
domain, and
\[
F_{w,v}(\xi)\;=\;\big(-\nabla_XU,\;+\nabla_YU\big),\qquad
S_{w,v}(\xi)\;=\;\Pi\big(\xi-\eta_xP_0F_{w,v}(\xi)\big)
\]
for the saddle operator of \eqref{eq:Udef} and the joint one-step location map.

\begin{lemma}[Curvature of the mixture objective]\label{lem:Ucurv}
Under Assumption~\ref{as:cells}(C1), on $\mathcal C$ the function $U$ of
\eqref{eq:Udef} is $m_\star$-strongly concave in $X$ and $m_\star$-strongly convex in
$Y$, with
\[
m_\star\;:=\;\min\{w_{\min},v_{\min}\}\cdot m\;\ge\;w_\star\,m\;>\;0,
\]
where $w_\star>0$ is the uniform lower bound on all weights from
Lemma~\ref{lem:full}(ii). Moreover $F_{w,v}$ is $L_F$-Lipschitz with $L_F\le L+L_{ab}$.
\end{lemma}
\begin{proof}
$U$ is separable across a player's own atoms, so $\nabla^2_{XX}U$ is the block diagonal
$\mathrm{diag}_k\big(w_k\nabla^2_{aa}u(x_k,\nu)\big)\preceq-w_{\min}mI$ by (C1), and
symmetrically $\nabla^2_{YY}U\succeq v_{\min}mI$. The mixed block has entries
$\nabla^2_{x_ky_l}U=w_kv_l\nabla^2_{ab}u(x_k,y_l)$; as a $K\times K$ array of blocks it
is dominated by the rank-one pattern $w\,v^\top$, so its operator norm is at most
$L_{ab}\norm{w}_2\norm{v}_2\le L_{ab}$. The diagonal blocks have norm at most
$\max_kw_kL\le L$.
\end{proof}

\begin{lemma}[The joint location step is a contraction]\label{lem:track}
Let $\eta_x\le m_\star p_{\min}/\big(p_{\max}^2L_F^2\big)$. Then
\begin{enumerate}[label=(\roman*),leftmargin=*]
\item \emph{(contraction)} $\norm{S_{w,v}(\xi)-S_{w,v}(\xi')}_{P_0^{-1}}
\le\big(1-\tfrac12\eta_x p_{\min}m_\star\big)\norm{\xi-\xi'}_{P_0^{-1}}$ for all
$\xi,\xi'\in\mathcal C$;
\item \emph{(perturbation)} for two weight pairs,
$\norm{S_{w,v}(\xi)-S_{w',v'}(\xi)}_{P_0^{-1}}\le\eta_x\sqrt{p_{\max}}\,
\norm{F_{w,v}(\xi)-F_{w',v'}(\xi)}_2$;
\item \emph{(fixed point)} $S_{w,v}$ has a unique fixed point in $\mathcal C$, at which
every $x_k$ maximises $u(\cdot,\nu)$ over $A_{c(k)}$ and every $y_l$ minimises
$u(\mu,\cdot)$ over $B_{c(l)}$.
\end{enumerate}
\end{lemma}
\begin{proof}
(i) The projection is nonexpansive, so drop it. Put $\zeta=\xi-\xi'$ and
$\Delta F=F_{w,v}(\xi)-F_{w,v}(\xi')$. Then
\[
\norm{\zeta-\eta_xP_0\Delta F}^2_{P_0^{-1}}
=\norm{\zeta}^2_{P_0^{-1}}-2\eta_x\ip{\zeta}{\Delta F}+\eta_x^2\norm{\Delta F}^2_{P_0}.
\]
The saddle operator of a function that is $m_\star$-strongly concave in $X$ and
$m_\star$-strongly convex in $Y$ is $m_\star$-strongly monotone: the two mixed terms
$-\ip{\nabla^2_{XY}U\,\zeta_Y}{\zeta_X}$ and $+\ip{\nabla^2_{YX}U\,\zeta_X}{\zeta_Y}$
enter with opposite signs and cancel, leaving only the diagonal blocks. Hence
$\ip{\zeta}{\Delta F}\ge m_\star\norm{\zeta}^2\ge m_\star p_{\min}\norm{\zeta}^2_{P_0^{-1}}$,
while $\norm{\Delta F}^2_{P_0}\le p_{\max}L_F^2\norm{\zeta}^2\le p_{\max}^2L_F^2\norm{\zeta}^2_{P_0^{-1}}$.
The squared factor is therefore at most
$1-2\eta_xm_\star p_{\min}+\eta_x^2p_{\max}^2L_F^2\le1-\eta_xm_\star p_{\min}$ under the
step-size condition, and $\sqrt{1-s}\le1-s/2$.

(ii) Drop the projection again:
$\norm{\eta_xP_0(F_{w,v}(\xi)-F_{w',v'}(\xi))}_{P_0^{-1}}
=\eta_x\norm{F_{w,v}(\xi)-F_{w',v'}(\xi)}_{P_0}\le\eta_x\sqrt{p_{\max}}\norm{\cdot}_2$.

(iii) Existence and uniqueness are immediate from (i) and compactness of $\mathcal C$.
At the fixed point, Lemma~\ref{lem:precond} gives coordinatewise stationarity of $U$,
which by \eqref{eq:Udef} and $w_k,v_l>0$ (Lemma~\ref{lem:full}(ii)) is stationarity of
$u(\cdot,\nu)$ at each $x_k$ and of $u(\mu,\cdot)$ at each $y_l$; (C1) and convexity of
the cells turn stationarity into the cell optimum.
\end{proof}

\begin{remark}[Where $L_{ab}$ went]\label{rem:noLab}
The cross-curvature $L_{ab}=\norm{\nabla^2_{ab}u}_\infty$ appears in
Lemma~\ref{lem:track} only inside $L_F$, i.e.\ only in the \emph{step-size} condition
$\eta_x\lesssim m_\star p_{\min}/(p_{\max}^2(L+L_{ab})^2)$. It does not appear in the
contraction factor and it imposes no threshold: for any $L_{ab}$, a small enough $\eta_x$
works. Analysing the two players' atoms separately --- bounding each one's motion as a
disturbance to the other's landscape --- instead produces a genuine constraint
$p_{\max}L_{ab}\lesssim p_{\min}m_\star$, unremovable by any step size. That constraint
is an artefact of the separate-ascent bookkeeping. The cancellation that removes it is a
game-theoretic fact about saddle problems, not an optimisation fact; see the second
bullet after Lemma~\ref{lem:identity}.
\end{remark}

\begin{remark}[Why the norm must be $\norm{\cdot}_{P_0^{-1}}$]\label{rem:Pnorm}
Statement (i) is false in the Euclidean norm for a general positive diagonal $P_0$:
$\ip{\zeta}{P_0\Delta F}$ can be bounded only by
$\big(p_{\min}m_\star-(p_{\max}-p_{\min})L_F\big)\norm{\zeta}^2$ from below, which is not
positive unless $P_0$ is nearly a multiple of the identity; and converting (i) back to
the Euclidean norm multiplies the contraction factor by $\sqrt{p_{\max}/p_{\min}}>1$,
destroying it. So the analysis stays in $\norm{\cdot}_{P_0^{-1}}$, which requires $P_0$
to be constant in $t$ --- a moving metric would reintroduce the same conversion factor at
every step. That is why \eqref{eq:Udef} matters: putting the $w_k$ in the gradient rather
than in the metric makes $P_0=\mathrm{diag}(\sigma_k^2)$, which is constant whenever
$\sigma$ is (Remark~\ref{rem:sigma}). The price is that $w_{\min}$ resurfaces in
$m_\star=w_{\min}m$ instead of in $p_{\min}$ --- the same product
$p_{\min}m_\star=\sigma_{\min}^2w_{\min}m$ either way.
\end{remark}

\newpage
\part*{Part III --- The theorem}
\addcontentsline{toc}{part}{Part III --- The theorem}

\section{Small gain: closing the loop}\label{sec:main}

We now do for the whole algorithm what Sokota et al.\ do for MMD on a fixed matrix:
exhibit \emph{one} fixed point and show that \emph{one} iteration moves the state
strictly closer to it. Nothing tracks a moving target, so neither the sensitivity of
the cell maximiser to $\nu$ nor the sensitivity of the QRE to $M$ is ever needed ---
and dropping those two sensitivity factors is what improves the final threshold from
$\kappa\lesssim m^2\tau^2$ to $\kappa\lesssim m\tau$ (Remark~\ref{rem:tau2}).

Write $T$ for one joint iteration of the algorithm, acting on
\[
\mathcal Z\;:=\;\mathcal C\times\simplex_K\times\simplex_K,
\qquad \mathcal C=\prod_{k\le K}A_{c(k)}\times\prod_{l\le K}B_{c(l)},
\]
a compact convex set by Assumption~\ref{as:cells} (cells are compact convex) and
Assumption~\ref{as:cover} (the assignment $c$ is well defined and $T$-invariant).

\begin{lemma}[A fixed point exists, and is what we want]\label{lem:exists}
$T$ is continuous and maps $\mathcal Z$ into $\mathcal Z$, so by Brouwer it has a fixed
point $\bar z=(\bar\xi,\bar w,\bar v)$, $\bar\xi=(\bar X,\bar Y)$. At any fixed point,
$\bar x_k=x^\star_{c(k)}(\bar\nu)$ and $\bar y_l=y^\star_{c(l)}(\bar\mu)$
(Lemma~\ref{lem:track}(iii)) and $(\bar w,\bar v)=\QRE_\tau\big(M(\bar X,\bar Y)\big)$.
\end{lemma}

\noindent Define the two error signals --- distances to \emph{that} fixed point, not to
anything moving:
\[
e_X(t):=\norm{\xi_t-\bar\xi}_{P_0^{-1}},
\qquad
e_W(t):=\sqrt{2\,D\big(\bar z_W\Vert z_W(t)\big)} .
\]
By Pinsker's inequality, $\norm{w_t-\bar w}_1+\norm{v_t-\bar v}_1\le e_W(t)$. Write
\[
G_\star\;:=\;\max\Big\{\max_k\norm{\nabla_au(\bar x_k,\bar\nu)},\;
\max_l\norm{\nabla_bu(\bar\mu,\bar y_l)}\Big\}
\]
for the residual gradient at the fixed point; $G_\star=0$ whenever every cell maximiser
is interior to $A$ (resp.\ $B$), which is the case of interest --- they are peaks --- and
$G_\star\le G$ always.

\begin{lemma}[Coupled recursion]\label{lem:coupled}
Under Assumptions~\ref{as:cells}--\ref{as:cover}, Proposition~\ref{prop:mmd}, and the
step sizes $\eta_x\le m_\star p_{\min}/(p_{\max}^2L_F^2)$ and $\eta_w\le\eta_0$,
\begin{align}
e_X(t+1) &\le \Big(1-\tfrac12\eta_x p_{\min}m_\star\Big)\,e_X(t)
\;+\;\eta_x\sqrt{p_{\max}}\,\Big(\tfrac{\kappa}{2}+G_\star\Big)\,e_W(t), \label{eq:recP}\\[2pt]
e_W(t+1) &\le \Big(1-\tfrac12 c\,\eta_w\tau\Big)\,e_W(t)
\;+\;8\,\eta_w\,G\sqrt{p_{\max}}\;e_X(t). \label{eq:recW}
\end{align}
\end{lemma}
\begin{proof}
\eqref{eq:recP}. Since $\bar z$ is a fixed point, $\bar\xi=S_{\bar w,\bar v}(\bar\xi)$,
while $\xi_{t+1}=S_{w_t,v_t}(\xi_t)$. By Lemma~\ref{lem:track}(i)--(ii),
\[
e_X(t+1)\;\le\;\big(1-\tfrac12\eta_xp_{\min}m_\star\big)e_X(t)
\;+\;\eta_x\sqrt{p_{\max}}\,\norm{F_{w_t,v_t}(\bar\xi)-F_{\bar w,\bar v}(\bar\xi)}_2 .
\]
For the $X$-block of that last operator difference, component $k$ is
\[
\bar w_k\nabla_au(\bar x_k,\bar\nu)-w_{t,k}\nabla_au(\bar x_k,\nu_t)
=(\bar w_k-w_{t,k})\,\nabla_au(\bar x_k,\bar\nu)
\;+\;w_{t,k}\big[\nabla_au(\bar x_k,\bar\nu)-\nabla_au(\bar x_k,\nu_t)\big].
\]
The first summand has norm at most $G_\star\abs{\bar w_k-w_{t,k}}$. For the second,
apply Lemma~\ref{lem:grad} with the \emph{same} atoms $\bar Y$ on both sides, so its
$L_{ab}$ term is identically zero and only $\tfrac{\kappa}{2}\norm{v_t-\bar v}_1$
remains. Stacking over $k$ and using $\norm{\cdot}_2\le\norm{\cdot}_1$ and
$\norm{w_t}_2\le1$,
\[
\norm{(F_{w_t,v_t}-F_{\bar w,\bar v})_X(\bar\xi)}_2
\le G_\star\norm{w_t-\bar w}_1+\tfrac{\kappa}{2}\norm{v_t-\bar v}_1 ,
\]
and symmetrically for the $Y$-block. Adding the two blocks and applying Pinsker gives
$\big(\tfrac{\kappa}{2}+G_\star\big)e_W(t)$.

\eqref{eq:recW}. Apply Lemma~\ref{lem:mmdpert} with $M=M(X_t,Y_t)$ and
$\bar M=M(\bar X,\bar Y)$, and use $\norm{z_W(t)-\bar z_W}_1\le e_W(t)$:
\[
\tfrac12e_W(t+1)^2\le(1-c\eta_w\tau)\tfrac12e_W(t)^2+2\eta_w\norm{M_t-\bar M}_{\max}e_W(t).
\]
If $\alpha^2\le(1-s)\beta^2+\gamma\beta$ with $s\in[0,1]$ then
$\alpha\le(1-\tfrac{s}{2})\beta+\gamma$, because
$\big((1-\tfrac{s}{2})\beta+\gamma\big)^2\ge(1-s)\beta^2+\gamma\beta$. Hence
$e_W(t+1)\le(1-\tfrac12c\eta_w\tau)e_W(t)+4\eta_w\norm{M_t-\bar M}_{\max}$. Finally
$\norm{M_t-\bar M}_{\max}\le G\big(\max_k\norm{x_k(t)-\bar x_k}+\max_l\norm{y_l(t)-\bar y_l}\big)
\le2G\sqrt{p_{\max}}\,e_X(t)$, since a coordinate max is at most the joint Euclidean
norm.
\end{proof}

\begin{theorem}[Convergence to the regularised fixed point]\label{thm:main}
Let Assumptions~\ref{as:cells} and \ref{as:cover} hold, let the step sizes satisfy the
conditions of Lemma~\ref{lem:coupled}, and suppose the \textbf{small-gain condition}
\begin{equation}\label{eq:gain}
\boxed{\;\mathcal G \;:=\; \frac{16\,(\kappa+2G_\star)\,G\,p_{\max}}{c\;p_{\min}\,m_\star\;\tau}\;<\;1 \;}
\end{equation}
holds. Then $\bar z$ is the \emph{unique} fixed point of $T$ in $\mathcal Z$, the iterates
converge to it linearly from every initialisation in $\mathcal Z$, at $\bar z$ every atom
sits at the optimum of its own cell against the opponent's final mixture and
$(\bar w,\bar v)=\QRE_\tau(M(\bar X,\bar Y))$, and
\[
\Expl(\bar\mu,\bar\nu)\;\le\;2\tau\log(1/\rho_{\min}).
\]
\end{theorem}
\begin{proof}
Lemma~\ref{lem:coupled} is $e(t+1)\le\mathcal M\,e(t)$ componentwise, with the
nonnegative matrix
$\mathcal M=\begin{pmatrix}1-\alpha_X & \beta_X\\ \beta_W & 1-\alpha_W\end{pmatrix}$,
$\alpha_X=\tfrac12\eta_xp_{\min}m_\star$,
$\beta_X=\eta_x\sqrt{p_{\max}}(\tfrac{\kappa}{2}+G_\star)$,
$\alpha_W=\tfrac12c\eta_w\tau$, $\beta_W=8\eta_wG\sqrt{p_{\max}}$. The step-size
conditions make both diagonal entries lie in $(0,1)$, so $I-\mathcal M$ is a $Z$-matrix
with positive diagonal, and $\rho(\mathcal M)<1$ iff its leading principal minors are
positive, i.e.\ iff $\det(I-\mathcal M)>0$, i.e.\ iff $\alpha_X\alpha_W>\beta_X\beta_W$:
\[
\tfrac12\eta_xp_{\min}m_\star\cdot\tfrac12c\eta_w\tau
\;>\;\eta_x\sqrt{p_{\max}}\big(\tfrac{\kappa}{2}+G_\star\big)\cdot 8\eta_wG\sqrt{p_{\max}}
\iff \tfrac14c\,p_{\min}m_\star\,\tau>4(\kappa+2G_\star)Gp_{\max}\iff\mathcal G<1 .
\]
\emph{Both step sizes cancel.} Hence $e(t)\le\mathcal M^te(0)\to0$ geometrically at rate
$\rho(\mathcal M)$, so $z_t\to\bar z$ linearly. Uniqueness is now free: if $\bar z'$ were
a second fixed point, the trajectory started at $\bar z'$ is constant and must converge
to $\bar z$, so $\bar z'=\bar z$. At $\bar z$: $e_X=0$ gives, by
Lemma~\ref{lem:exists} and \eqref{eq:cellmax}, $\Delta_x=\Delta_y=0$; $e_W=0$ gives
$(\bar w,\bar v)=\QRE_\tau$, so $\Expl_M\le2\tau\log(1/\rho_{\min})$ by
Lemma~\ref{lem:bias}. Conclude with the master identity, Lemma~\ref{lem:identity}.
\end{proof}

\begin{remark}[Read \eqref{eq:gain} as $\kappa \lesssim m_\star\,\tau$]\label{rem:tau2}
Take $G_\star=0$ (interior cell maxima) and treat $p_{\max}/p_{\min}$, $G$ and $c$ as
$O(1)$. The small-gain condition reads
\[
\kappa \;\lesssim\; m_\star\,\tau,\qquad m_\star=w_{\min}\,m .
\]
\textbf{Coupling must be small compared with curvature $\times$ temperature.} That is the
entire game-theoretic content of the theorem. Note that $m_\star$, not $m$, is the
relevant curvature: low-weight atoms move slowly, and that shows up in the
\emph{stability threshold}, not merely in the rate --- one more reason for the change
recommended in \S\ref{sec:starve}, which replaces $m_\star$ by $m$ outright.

An earlier version of this note ran the same argument through moving targets: it
measured $e_X$ against the \emph{current} cell maximiser $x^\star_{c(k)}(\nu_t)$ and
$e_W$ against the \emph{current} $\QRE_\tau(M_t)$, and therefore had to pay for how fast
each target moves --- a factor $\kappa/m$ for the cell maximiser and a factor $1/\tau$ for
the QRE. Those two factors turned the threshold into $\kappa\lesssim m^2\tau^2$. They are
an artefact of the bookkeeping, not of the algorithm: measuring both errors against the
fixed point of the joint map removes both, at the cost of needing Brouwer for existence
(Lemma~\ref{lem:exists}) and of the perturbed-step Lemma~\ref{lem:mmdpert} in place of a
QRE sensitivity bound. Since $m_\star\tau<1$ in every regime of interest,
$\kappa\lesssim m_\star\tau$ is the strictly weaker, i.e.\ better, requirement.
\end{remark}

\begin{remark}[Only one threshold, and it involves only $\kappa$]
\eqref{eq:gain} is now the sole non-step-size condition of the theorem. In particular
the cross-curvature $L_{ab}$ does \emph{not} appear: by Remark~\ref{rem:noLab} it lives
in the step-size condition through $L_F\le L+L_{ab}$ and nowhere else. This is worth
stating because the opposite is easy to believe --- and an earlier draft did assert a
second threshold $p_{\max}L_{ab}\le\tfrac14p_{\min}m$. What kills it is that the two
players' location blocks are ascent and descent on the \emph{same} function
\eqref{eq:Udef}, so their mutual influence is antisymmetric and cancels in
$\ip{F(\xi)-F(\xi')}{\xi-\xi'}$. Treating each side as an exogenous disturbance to the
other throws that cancellation away.
\end{remark}

\begin{remark}[The step sizes cancel --- timescale separation does not save you]
This is the least obvious consequence of Lemma~\ref{lem:coupled} and worth stating
loudly. One's instinct is that a two-timescale scheme ($\eta_x\ll\eta_w$, or the
reverse) should make the interaction negligible. It does not, because both the
damping and the disturbance in each loop scale with the \emph{same} step size:
slowing the atoms down slows their disturbance of the weights by exactly as much as
it slows their own recovery. The loop gain $\mathcal G$ is step-size free. Timescale
separation changes the transient, the constants, and the rate --- never the
threshold. What \emph{does} move the threshold is $\tau$ (raise it), the curvature
$m_\star$ (a property of the landscape and of the weights; see \S\ref{sec:anneal-sigma}
and \S\ref{sec:starve}), and the coupling $\kappa$ (a property of the game).
\end{remark}

\section{From QRE to Nash}\label{sec:anneal-tau}

Theorem~\ref{thm:main} gives a $2\tau\log(1/\rho_{\min})$-Nash. Two routes to $0$.

\paragraph{Route 1: magnet refreshes at fixed $\tau$.} This is what the code does
(\texttt{magnet\_interval}). Setting $\rho\leftarrow w$ every $T_{\mathrm{mag}}$
steps makes the outer iteration a proximal-point method on the restricted game, which
converges to its exact Nash --- provided $T_{\mathrm{mag}}$ is long enough that the
inner loop has converged, which by Theorem~\ref{thm:main} takes
$O(\log(1/\eps)/\alpha)$ steps. The catch is Lemma~\ref{lem:full}: the bound
$\hat w_k\ge\rho_k e^{-2U_{\max}/\tau}$ degrades geometrically as $\rho$ itself
approaches a Nash with zero-weight atoms.

\paragraph{Route 2: anneal $\tau_t\downarrow0$ with a uniform magnet.} Cleaner to
analyse, and it exposes the real tension.

\begin{proposition}[Annealing schedule]\label{prop:anneal}
Suppose $\tau_t\downarrow0$ with $\tau_t$ non-increasing and
$\tau_t-\tau_{t+1}\le \eps_\tau\,\tau_t$. All the constants of
Lemma~\ref{lem:coupled} hold with $\tau=\tau_t$, plus one extra disturbance term
of size $O\big((\tau_t-\tau_{t+1})/\tau_t\big)$ in \eqref{eq:recW}. Two conditions
must hold along the schedule:
\begin{enumerate}[label=(\roman*),leftmargin=*]
\item \emph{stability}: $\kappa\lesssim m_\star\tau_t$ --- so $\tau_t$ may not fall
below $\tau_\infty\asymp\kappa/m_\star$ unless the game's coupling is exactly zero;
\item \emph{no starvation}: the atoms must keep up, and their speed is
$\propto p_{\min}=\sigma^2_{\min}w_{\min}$ with $w_{\min}\ge K^{-1}e^{-2U_{\max}/\tau_t}$
by Lemma~\ref{lem:full}, so the target's drift per step must satisfy
$\eps_\tau \lesssim \eta_x\,e^{-2U_{\max}/\tau_t}$.
\end{enumerate}
Condition (ii) alone forces $\tau_t\gtrsim 2U_{\max}/\log t$, giving
$\Expl_t = O\big(U_{\max}\log K/\log t\big)$. Note that (ii) only bites in the regime
that (i) permits: if $\kappa>0$ then (i) stops the schedule at $\tau_\infty\asymp\kappa/m_\star$
long before the $1/\log t$ rate becomes the binding constraint, and (ii) is the whole
story only for $\kappa=0$ (or after the fix in \S\ref{sec:starve}, which removes it).
\end{proposition}

\noindent Both conditions are informative and both are actionable.

\begin{itemize}[leftmargin=*]
\item \textbf{(i) is a real floor.} If the game has genuine coupling
($\kappa>0$) and finite curvature, this analysis does not give exact Nash --- it gives
$\Expl = O(\kappa\log K/m_\star)$. This is not obviously an artefact: it is the same
place where \texttt{theory2} needs its stability condition, and it is where the
counterexample games in \texttt{configs/idealized\_counterexample/} live.
\item \textbf{(ii) is an artefact of the preconditioner, and is fixable.} The
$e^{-2U_{\max}/\tau}$ comes purely from the factor $w_k$ in \eqref{eq:xstep}:
low-weight atoms crawl. \textbf{Normalising the per-component location gradient by
$w_k$} --- i.e. taking $p_k=\sigma_k^2$ instead of $p_k=\sigma_k^2 w_k$ --- removes
it entirely, is legitimate by Lemma~\ref{lem:precond} (it is a different positive
preconditioner, so the fixed points are identical), and upgrades
Proposition~\ref{prop:anneal} from a $1/\log t$ rate to a polynomial one. See
\S\ref{sec:starve}.
\end{itemize}

\newpage
\part*{Part IV --- What breaks it, and what it says about the code}
\addcontentsline{toc}{part}{Part IV --- What breaks it, and what it says about the code}

\section{Failure modes, each one an assumption being violated}

\subsection{Two maxima in one cell (Assumption~\ref{as:cells} (C1) fails)}
If a region that the flow cannot leave contains two local maxima, an atom
initialised in the wrong basin sits at a strictly suboptimal stationary point and
$\Delta_x$ does not close. Note what Lemma~\ref{lem:identity} says about this: the
weight loop \emph{cannot compensate}. $\Expl_M$ may be exactly $0$ --- the weights
are a perfect Nash of the restricted game --- while $\Delta_x>0$ forever. Adding
more atoms does not help unless one of them lands in the good basin, which is a
statement about initialisation, not about $K$. This is the ``decoy'' family
(\texttt{DecoyWellGame}, \texttt{configs/idealized\_decoy\_well.yaml}).

\subsection{Coupling stronger than curvature (\eqref{eq:gain} fails)}
When $\kappa \gtrsim m_\star\tau$ the cell maximisers $x_i^\star(\nu)$ swing further
than the atoms can track, the restricted game $M(X_t,Y_t)$ never settles, and the
two loops chase each other. This is the classical matching-pennies rotation, lifted
into the locations. Symptom in a run: exploitability plateaus with visible
oscillation in \texttt{means0}/\texttt{means1} rather than decaying.

\subsection{Flat landscapes ($m=0$)}
If $u(\cdot,\nu)$ has no curvature there is no cell maximiser to track and
Lemma~\ref{lem:track} is vacuous. The atoms drift on a plateau, driven only by the
coupling term. Note $\mathcal G\to\infty$ as $m\to0$: the theory says flat games are
the \emph{worst} case, not the easiest. This matches the empirical observation that
the dead plateau between two narrow wells is where the algorithm gets stuck.

\subsection{Weight starvation ($p_{\min}\to0$)}\label{sec:starve}
The location step \eqref{eq:xstep} carries a factor $w_k$. An atom whose weight
collapses stops moving, so it stops improving, so its weight stays collapsed --- a
self-reinforcing trap. Lemma~\ref{lem:full} says the magnet/entropy temperature
$\tau$ is precisely what prevents this ($w_k\ge\rho_ke^{-2U_{\max}/\tau}$), which is
a second, independent reason to keep $\tau$ away from $0$ beyond the stability
condition \eqref{eq:gain}.

\medskip
\noindent\textbf{Concrete recommendation for the codebase.} In
\texttt{gaussian\_natural\_step}, the mean gradient is
$w_k\sigma_k^2\,\nabla_a u(x_k,\nu)$. Dividing by $w_k$ (equivalently, computing each
component's location gradient against its \emph{own} conditional objective rather
than the mixture's) is free by Lemma~\ref{lem:precond} and removes both the
$e^{-2U_{\max}/\tau}$ in Proposition~\ref{prop:anneal} and the entire starvation
failure mode. It also replaces $m_\star=w_{\min}m$ by $m$ in the effective curvature
of Lemma~\ref{lem:Ucurv}: with $p_k=\sigma_k^2$ the location block ascends
$\sum_ku(x_k,\nu)$, on which each atom feels the full curvature $m$ regardless of its
weight. So it moves the \emph{stability threshold} \eqref{eq:gain}, not just the rate ---
by a factor $1/w_{\min}$, which Lemma~\ref{lem:full} bounds only by
$Ke^{2U_{\max}/\tau}$. That is the largest single constant in this theory, and it is
free to remove. It should show up as: low-weight atoms still migrating to their cell
maxima, so that a later change in the opponent can revive them.

\subsection{Too few atoms, or a bad initialisation (Assumption~\ref{as:cover} fails)}
Ascent never moves an atom across a cell boundary, so an uncovered cell is
uncovered forever, and if that cell holds the global best response then
$\Delta_x>0$ permanently. Everything about ``more points than you need'' is used
\emph{here} and nowhere else in the proof.

\section{What the Gaussian width actually does}\label{sec:anneal-sigma}

By Remark~\ref{rem:sigma} all constants are those of $u_\sigma$, and the width
changes them in a specific, useful way. For a landscape built from Gaussian bumps of
width $\omega$ --- exactly \texttt{MultiPointGame} --- smoothing by $\sigma$ gives
bumps of width $\sqrt{\omega^2+\sigma^2}$ and height scaled by
$\omega/\sqrt{\omega^2+\sigma^2}$ (this is \texttt{well\_expectation}). Therefore:
\[
\text{cell radius } \approx \sqrt{\omega^2+\sigma^2},\qquad
m \;\approx\; \frac{h\,\omega}{(\omega^2+\sigma^2)^{3/2}},\qquad
m_A(\sigma)\ \text{non-increasing in }\sigma .
\]
So:
\begin{itemize}[leftmargin=*]
\item \textbf{Large $\sigma$: few, wide cells.} Assumption~\ref{as:cover} is easy
(a handful of atoms covers everything), and there is no dead plateau --- an atom
anywhere feels the far peak. But nearby peaks have merged into one cell, so the
$\sigma$-landscape's maxima are not the true maxima, and the fixed point is
biased.
\item \textbf{Small $\sigma$: many, narrow cells, high curvature.} $m$ is large, so
\eqref{eq:gain} is easy to satisfy and the fixed point is unbiased --- but
Assumption~\ref{as:cover} needs an atom inside each narrow well, and between the
wells the landscape is flat, so an atom that starts outside every well never
arrives.
\end{itemize}
This is exactly graduated optimisation, and it is what \texttt{anneal\_std\_from}
implements: start with wide cells so that coverage is automatic, then shrink,
splitting cells while each already contains an atom.

\begin{remark}[The honest status of $\sigma$-annealing]
Theorem~\ref{thm:main} assumes the cell decomposition is \emph{fixed}. A
$\sigma$-schedule changes it, so the theorem does not cover annealed runs. The
natural extension --- and, I think, the right next thing to prove --- is a
\emph{cell-splitting} statement:
\begin{quote}
if $\sigma$ decreases slowly enough that, at each splitting of a cell $A_i$ into
$A_{i'},A_{i''}$, at least one atom lies in each part, and the atoms have converged
to within $O(\text{cell radius})$ before each split, then coverage
(Assumption~\ref{as:cover}) is preserved for all $\sigma$ and
Theorem~\ref{thm:main} applies at every stage.
\end{quote}
This requires $K\ge m_A(\sigma_{\min})$ (again ``more points than you need''), plus a
mechanism to leave more than one atom per cell during the wide phase --- which is
what the magnet's full-support guarantee (Lemma~\ref{lem:full}) and any repulsion
between components would provide. This is stated as a conjecture, not proved.
\end{remark}

\section{Checking the assumptions on this codebase's game}

\texttt{MultiPointGame} with peaks $p_1<\dots<p_n$, bump width $\omega$, coupling
$\lambda$:
\[
u(a,b) \;=\; D(a)-D(b)+\lambda\,f(a)^\top f(b),\qquad
D(a)=\sum_{j}e^{-(a-p_j)^2/2\omega^2},\quad f_j(a)=\Big(\tfrac{a-\bar p}{r}\Big)^{j}-t_j .
\]
Against a fixed $\nu$, the maximiser's landscape is
$a\mapsto D(a)+\lambda f(a)^\top\bar f_\nu + \text{const}$ with
$\bar f_\nu=\E_\nu[f]$, $\norm{\bar f_\nu}_\infty\le2$.

\paragraph{(C1) curvature.} $D''(p_j)=-1/\omega^2$, and $D$ is strongly concave with
constant $\asymp1/\omega^2$ on $\lvert a-p_j\rvert\le\omega$. The coupling term
contributes at most $\lambda\,\norm{f''}_\infty\norm{\bar f_\nu}_\infty
\le 2\lambda\,C_n/r^2$ where $r$ is the half-range and $C_n$ depends on the polynomial
degree. So (C1) holds on $A_j=[p_j-\omega,p_j+\omega]$ with
$m \asymp \omega^{-2}-2\lambda C_n/r^2$, i.e. whenever
\begin{equation}\label{eq:mpgcond}
\lambda \;\lesssim\; \frac{r^2}{\omega^2}\quad\text{(up to the degree constant }C_n).
\end{equation}
For the defaults (peaks $\pm1$ so $r=1$, $\omega=0.1$, $\lambda=1$) this is
comfortable: $r^2/\omega^2=100$.

\paragraph{(C2)/(C3) invariance and the plateau.} The cells $[p_j\pm\omega]$ do not
cover $A$: between and outside the wells $D$ is exponentially flat, and only the
coupling term $\lambda f'(a)^\top\bar f_\nu$ pushes. Per Remark~\ref{rem:nocover} this
is not a defect of the assumption --- coverage would contradict (C1) --- but it does
mean (C3) has real content here: the global maximiser must sit inside some
$[p_j\pm\omega]$, which holds as long as the coupling term does not create a maximum on
the plateau, i.e.\ under \eqref{eq:mpgcond}. What is genuinely not supplied by the game
is Assumption~\ref{as:cover}: an atom that starts on the plateau belongs to no cell and
the theory says nothing about it. Two ways to fix that, both already in the code: take
$\sigma\gtrsim$ (peak separation)$/3$ so that $u_\sigma$'s wells overlap and the
cells tile $A$ (\S\ref{sec:anneal-sigma}); or initialise atoms inside the wells
(\texttt{\_spread\_bias\_init} at the peaks), which is Assumption~\ref{as:cover}
by construction.

\paragraph{The small-gain number.} $\kappa \le 2\lambda\norm{f'}_\infty\norm{f}_\infty
\asymp\lambda/r$, so \eqref{eq:gain} reads $\kappa\lesssim m_\star\tau$ with
$m_\star=w_{\min}m$, i.e.
\[
\frac{\lambda}{r} \;\lesssim\; \frac{w_{\min}}{\omega_\sigma^2}\,\tau
\quad\Longleftrightarrow\quad
\tau \;\gtrsim\; \frac{\lambda\,\omega_\sigma^2}{r\,w_{\min}},
\qquad \omega_\sigma:=\sqrt{\omega^2+\sigma^2},
\]
where $\omega_\sigma$ rather than $\omega$ appears because all constants are those of
$u_\sigma$ (Remark~\ref{rem:sigma}). Three numbers go in. With $\omega=0.1$ and the
code's typical $\sigma\approx0.1$, $\omega_\sigma^2\approx0.02$. With $\lambda=r=1$ the
coupling is $\kappa\asymp1$. For $w_{\min}$, the a priori bound of
Lemma~\ref{lem:full} ($K^{-1}e^{-2U_{\max}/\tau}$) is numerically vacuous here, so use
the value the run actually sits at: with $K=2$ and a near-uniform equilibrium mixture,
$w_{\min}\approx\tfrac12$. That gives
\[
\tau \;\gtrsim\; 0.04 .
\]
The code's default \texttt{magnet\_coef} is $0.2$, a factor of $5$ above the threshold
--- which predicts that the default configuration is in the stable regime, and that
pushing $\tau$ down through $\sim\!4\times10^{-2}$ should destabilise it. \emph{That is a
falsifiable prediction and a cheap experiment: sweep \texttt{magnet\_coef} down and look
for the onset of oscillation in exploitability.} The margin is thinner than it looks in
Remark~\ref{rem:tau2}, and the reason is entirely the $w_{\min}$ in $m_\star$: with the
\S\ref{sec:starve} change ($p_k=\sigma_k^2$) the prediction moves back to
$\tau\gtrsim0.02$ and stops depending on the equilibrium mixture at all. Comparing the
two sweeps is a second, sharper experiment.

\section{What this theory deliberately does not prove}

\begin{enumerate}[leftmargin=*]
\item \textbf{Exact Nash with $\kappa>0$.} Proposition~\ref{prop:anneal}(i) leaves a
floor $\Expl=O(\kappa\log K/m_\star)$. Removing it needs either an argument that
the small-gain bound is loose (it is: $\mathcal G$ uses worst-case $\kappa$ over all
$\nu$, whereas near the fixed point the relevant $\nu$-variation is small) or a
genuine last-iterate result for the coupled system.
\item \textbf{Sampling and approximation error.} Assumed away by construction, per
the brief. Adding them is a routine stochastic-approximation exercise on top of
Lemma~\ref{lem:coupled}: both recursions become $\le(1-\alpha)e+\text{disturbance}+\text{noise}$
and one gets an $O(\text{noise}/\alpha)$ radius.
\item \textbf{$\sigma$-annealing.} Conjecture only, \S\ref{sec:anneal-sigma}.
\item \textbf{Non-covering initialisations, and any escape mechanism.} Ascent cannot
cross cells; nothing here provides exploration between cells, and
Assumption~\ref{as:cover} simply asserts it away.
\item \textbf{Extensive form / states.} Single-shot games only. The master identity
(Lemma~\ref{lem:identity}) is stated per information set and should compose, but
that is not done here.
\end{enumerate}

\section{The whole thing on one page}

\begin{enumerate}[leftmargin=*]
\item \textbf{Identity.} $\Expl = \Expl_{M(X,Y)}(w,v) + \Delta_x + \Delta_y$, exactly.
Mixing error and support error add; there is no cross term
(Lemma~\ref{lem:identity}).
\item \textbf{Weights kill the first term.} MMD with temperature $\tau>0$ is a
contraction to $\QRE_\tau$ of the restricted $K\times K$ game, with bias
$2\tau\log K$ (Proposition~\ref{prop:mmd}, Lemma~\ref{lem:bias}). The magnet is not
a heuristic --- without it there is no contraction at all.
\item \textbf{Locations kill the other two.} Each atom does preconditioned ascent on
$u(\cdot,\nu)$, which inside a cell is a contraction --- in the preconditioner's own
norm, Remark~\ref{rem:Pnorm} --- to that cell's maximiser (Lemma~\ref{lem:track}).
Since the global max is attained in some cell, covering the cells makes $\Delta_x=0$
(Assumptions~\ref{as:cells}--\ref{as:cover}). \emph{This is
the only place the ``more points than you need'' hypothesis is used.}
\item \textbf{The Gaussian head is a preconditioner.} It multiplies the gradient by
$\sigma_k^2w_k>0$, which cannot move a fixed point (Lemma~\ref{lem:precond}). It
enters the result only through $p_{\max}/p_{\min}$, and --- separately --- through
which smoothed landscape the cells are defined on.
\item \textbf{Coupling of the two loops.} Small gain:
$\kappa \lesssim m_\star\tau$, $m_\star=w_{\min}m$ (Theorem~\ref{thm:main},
Remark~\ref{rem:tau2}) --- the \emph{only} non-step-size condition.
Coupling versus curvature, moderated by temperature. \textbf{Both step sizes cancel:
timescale separation does not change the threshold.} The whole coupling argument is one
contraction towards one fixed point, exactly as in the fixed-matrix MMD proof; nothing
tracks a moving target.
\item \textbf{Price.} $\Expl \le 2\tau\log K$ at the fixed point; annealing $\tau$
is limited by stability ($\tau\gtrsim\kappa/m_\star$) and by weight starvation, the
latter removable by dropping $w_k$ from the location preconditioner.
\end{enumerate}

\end{document}
